\section*{Bab \ref{ch:g6:cell-unit-of-life} --- Sel, Satuan Kehidupan}

\begin{solution}{exo:g6:cell-unit-of-life:1}
\omterm{def:g6:cell-unit-of-life:parts}{Membran} --- selubung tipis yang menahan \omterm{def:g5:first-look-at-cells:cell}{selnya} tetap menyatu dan
mengendalikan apa yang masuk serta keluar; \omterm{def:g6:cell-unit-of-life:parts}{sitoplasma} --- isian mirip
agar-agar tempat pekerjaannya berlangsung; \omterm{def:g6:cell-unit-of-life:parts}{inti sel} --- pusat komando
yang mengarahkan kehidupan \omterm{def:g5:first-look-at-cells:cell}{selnya} dan pembelahannya.
\end{solution}

\begin{solution}{exo:g6:cell-unit-of-life:2}
\omterm{prop:g6:cell-unit-of-life:plantextras}{Dinding sel} yang kaku --- kekakuan \omterm{def:g1:plants-around-us:plant}{tumbuhan} tanpa \omterm{def:g3:bones-and-muscles:skeleton}{tulang}; butir hijau
--- bengkel penangkap \omterm{def:g6:exploring-our-environment:conditions}{cahaya} bagi perdagangan para \omterm{def:g5:food-webs:roles}{produsen}; kantong
air yang besar --- tekanannya menjaga \omterm{def:g5:first-look-at-cells:cell}{selnya} tetap gemuk.
\end{solution}

\begin{solution}{exo:g6:cell-unit-of-life:3}
Bawang bombai: dinding lurus yang bertemu pada sudutnya, \omterm{def:g6:cell-unit-of-life:parts}{inti selnya}
terdesak ke tepi. Pipi: garis luar yang membulat dan lunak, hanya
bermembran, \omterm{def:g6:cell-unit-of-life:parts}{inti selnya} di dekat tengahnya.
\end{solution}

\begin{solution}{exo:g6:cell-unit-of-life:4}
Sebuah organisme lengkap yang terdiri atas persis satu \omterm{def:g5:first-look-at-cells:cell}{sel}. \omterm{ex:g6:producing-our-food:bread}{Raginya};
si pendayung berbentuk sandal di setetes air kolamnya (atau perenang
hijaunya, gumpalan yang mengalir, \omterm{def:g6:producing-our-food:fermentation}{mikroba} yogurtnya).
\end{solution}

\begin{solution}{exo:g6:cell-unit-of-life:5}
Setiap \omterm{def:g6:exploring-our-environment:components}{makhluk hidup} tersusun atas satu \omterm{def:g5:first-look-at-cells:cell}{sel} atau lebih, yaitu satuan
kehidupan yang terkecil; dan setiap \omterm{def:g5:first-look-at-cells:cell}{sel} datang dari pembelahan sebuah
\omterm{def:g5:first-look-at-cells:cell}{sel} yang sudah ada.
\end{solution}

\begin{solution}{exo:g6:cell-unit-of-life:6}
Dengan membelah menjadi dua: satu \omterm{def:g5:first-look-at-cells:cell}{sel} menjadi dua organisme yang
lengkap.
\end{solution}

\begin{solution}{exo:g6:cell-unit-of-life:7}
\omterm{prop:g6:cell-unit-of-life:plantextras}{Dinding sel} yang kaku dan kantong air yang gemuk: jutaan \omterm{def:g5:first-look-at-cells:cell}{sel} yang
berdinding dan bertekanan saling menopang seperti susunan bata yang
digelembungkan. Layu adalah kantongnya yang mengendur --- tekanannya
hilang karena kekurangan air, dan seluruh bangunannya melorot.
\end{solution}

\begin{solution}{exo:g6:cell-unit-of-life:8}
Butir hijaunya. \omterm{def:g1:animals-around-us:animal}{Hewan} adalah \omterm{def:g5:food-webs:roles}{konsumen} --- pengubah \omterm{def:g6:plants-are-producers:matter}{materi organik} yang
dimakan, tidak pernah pembuat-dari-mineral --- sehingga \omterm{def:g5:first-look-at-cells:cell}{selnya} tidak
membawa butir pembuat: perlengkapannya cocok dengan perdagangannya,
\omterm{def:g5:first-look-at-cells:cell}{sel} demi \omterm{def:g5:first-look-at-cells:cell}{sel}.
\end{solution}

\begin{solution}{exo:g6:cell-unit-of-life:9}
Misalnya: \omterm{def:g5:first-look-at-cells:cell}{sel} pelapis (\omterm{def:g5:first-look-at-cells:cell}{sel} pipinya), yang menutupi permukaan; \omterm{def:g5:first-look-at-cells:cell}{sel}
\omterm{def:g3:bones-and-muscles:muscle}{otot}, yang \omterm{def:g3:bones-and-muscles:muscle}{berkontraksi} untuk menarik \omterm{def:g3:bones-and-muscles:skeleton}{tulang}; \omterm{def:g5:first-look-at-cells:cell}{sel} \omterm{def:g5:brain-in-command:system}{saraf}, yang membawa
laporan dan perintah pada gelang \omterm{def:g5:brain-in-command:system}{otaknya}.
\end{solution}

\begin{solution}{exo:g6:cell-unit-of-life:10}
Lecet yang sembuh: \omterm{def:g5:first-look-at-cells:cell}{sel} \omterm{def:g1:the-five-senses:sense}{kulit} yang membelah untuk membangun ulang
petaknya. Adonan yang berlipat dua: \omterm{def:g5:first-look-at-cells:cell}{sel} \omterm{ex:g6:producing-our-food:bread}{ragi} yang membelah lagi dan
lagi di dalam adonan hangat yang diberi gula. \omterm{ex:g2:animals-grow-up:frog}{Berudu} yang tumbuh:
\omterm{def:g5:first-look-at-cells:cell}{selnya} yang membelah lalu mengkhusus sementara tubuhnya dibangun.
\end{solution}

\begin{solution}{exo:g6:cell-unit-of-life:11}
Perdagangan para \omterm{def:g5:food-webs:roles}{produsen}: pembuatan \omterm{def:g6:plants-are-producers:matter}{materi organik} dari pasokan
mineral dan \omterm{def:g6:exploring-our-environment:conditions}{cahaya} oleh \omterm{def:g5:first-look-at-cells:cell}{sel} tunggal. Di dalam jaring kolamnya mereka
adalah mata rantai pertama --- lantai hijau yang dirumput \omterm{ex:g2:animals-grow-up:frog}{berudu} dan
kutu air.
\end{solution}

\begin{solution}{exo:g6:cell-unit-of-life:12}
Tiap selmu datang dari sebuah pembelahan, yang \omterm{def:g5:first-look-at-cells:cell}{sel} induknya datang
dari sebuah pembelahan pula, tanpa satu pun putus --- mundur melewati
\omterm{def:g4:animal-reproduction:sexual}{sel telur} yang terbuahi, melewati \omterm{def:g5:first-look-at-cells:cell}{sel} orang tuamu, dan lebih jauh lagi
ke masa lalu kehidupan yang dalam. Akibatnya dalam satu kalimat:
setiap \omterm{def:g5:first-look-at-cells:cell}{sel} yang hidup, termasuk selmu, adalah mata rantai terbaru dari
satu rantai pembelahan yang tak terputus, yang setua kehidupan itu
sendiri.
\end{solution}

\begin{solution}{pb:g6:cell-unit-of-life:1}
\textbf{1.} Dindingnya: \omterm{prop:g6:cell-unit-of-life:plantextras}{dinding sel} \omterm{def:g1:plants-around-us:plant}{tumbuhannya}. Bintiknya: \omterm{def:g6:cell-unit-of-life:parts}{inti
selnya}. Pelaku pendesakannya: kantong air yang besar, yang menekan
\omterm{def:g6:cell-unit-of-life:parts}{inti selnya} ke tepi.

\textbf{2.} A punya dinding, kantong air (dan, pada bagian yang hijau,
butir hijau) yang tidak dimiliki B. Keduanya berbagi perlengkapannya:
\omterm{def:g6:cell-unit-of-life:parts}{membran}, \omterm{def:g6:cell-unit-of-life:parts}{sitoplasma}, \omterm{def:g6:cell-unit-of-life:parts}{inti sel}.

\textbf{3.} Butirnya adalah bengkel penangkap \omterm{def:g6:exploring-our-environment:conditions}{cahaya} --- dan sebuah
umbi bawah \omterm{def:g6:decomposers-and-soil:soil}{tanah} hidup dalam gelap, tempat bengkelnya akan menganggur:
\omterm{def:g1:plants-around-us:plant}{tumbuhannya} membangun bengkelnya di \omterm{def:g1:plants-around-us:parts}{daun} yang tersinari, bukan di
gudangnya yang terkubur.

\textbf{4.} \omterm{def:g5:first-look-at-cells:cell}{Sel} murid-muridnya sendiri. Menurut kalimat 2,
masing-masingnya datang dari pembelahan sebuah \omterm{def:g5:first-look-at-cells:cell}{sel} yang sudah ada ---
sebuah rantai pembelahan yang menjangkau kembali sampai ke \omterm{def:g4:animal-reproduction:sexual}{sel telur}
terbuahi milik tiap muridnya sendiri.

\textbf{5.} Si pendayung: bergerak, mengemudi, makan --- sebuah
organisme \omterm{def:g6:cell-unit-of-life:unicellular}{bersel tunggal} yang lengkap (yang mirip \omterm{def:g1:animals-around-us:animal}{hewan}). Para
perenang hijau: bergerak dan menangkap \omterm{def:g6:exploring-our-environment:conditions}{cahaya} --- \omterm{def:g5:food-webs:roles}{produsen} \omterm{def:g6:cell-unit-of-life:unicellular}{bersel
tunggal}. Gumpalannya: mengalir, menelan makanan --- seorang pemburu
\omterm{def:g6:cell-unit-of-life:unicellular}{bersel tunggal}. Masing-masingnya memperlihatkan kegiatan sebuah
kehidupan yang utuh di dalam satu \omterm{def:g5:first-look-at-cells:cell}{sel}.

\textbf{6.} Butir hijaunya --- perlengkapan penangkap \omterm{def:g6:exploring-our-environment:conditions}{cahaya} yang
menjalankan pembuatan oleh para \omterm{def:g5:food-webs:roles}{produsen}.

\textbf{7.} Mereka makan dari gula \omterm{def:g2:animals-grow-up:born}{susunya} lalu melepaskan asam yang
memadatkannya menjadi yogurt: mereka adalah \omterm{def:g6:exploring-our-environment:components}{makhluk hidup} \omterm{def:g6:cell-unit-of-life:unicellular}{bersel
tunggal} --- angkatan kerja \omterm{def:g6:producing-our-food:fermentation}{mikrobanya}, yang akhirnya terlihat langsung.

\textbf{8.} Amatilah sebuah bintik membelah: satu menjadi dua, dua
menjadi empat --- pembelahan adalah tanda tangan yang dituntut kalimat
2 dari yang hidup. (Di atas kaca preparat yang hangat, \omterm{def:g6:producing-our-food:fermentation}{mikroba}
yogurtnya menurut dalam hitungan jam.)

\textbf{9.} A: banyak \omterm{def:g5:first-look-at-cells:cell}{sel}, \omterm{def:g1:plants-around-us:plant}{tumbuhan}. B: banyak \omterm{def:g5:first-look-at-cells:cell}{sel} (sebuah contoh
\omterm{def:g1:animals-around-us:animal}{hewan} bersel banyak), \omterm{def:g1:animals-around-us:animal}{hewan}. C: \omterm{def:g5:first-look-at-cells:cell}{sel} tunggal --- yang mirip \omterm{def:g1:animals-around-us:animal}{hewan} (si
pendayung, gumpalannya) dan yang mirip \omterm{def:g1:plants-around-us:plant}{tumbuhan} (para perenang hijau)
berdampingan. D: \omterm{def:g5:first-look-at-cells:cell}{sel} tunggal, bukan \omterm{def:g1:plants-around-us:plant}{tumbuhan} dan bukan \omterm{def:g1:animals-around-us:animal}{hewan} ---
\omterm{def:g6:producing-our-food:fermentation}{mikroba}.

\textbf{10.} \omterm{def:g6:cell-unit-of-life:parts}{Membran}, \omterm{def:g6:cell-unit-of-life:parts}{sitoplasma}, \omterm{def:g6:cell-unit-of-life:parts}{inti sel}. Kesemestaannya melintasi
bawang bombai, pipi, kolam dan yogurt adalah bukti \omterm{prop:g5:first-look-at-cells:principle}{kesatuan kehidupan}:
satu rancangan pembangunan di bawah semuanya.

\textbf{11.} Bawang bombainya tumbuh dari satu \omterm{def:g5:first-look-at-cells:cell}{sel} buatan \omterm{def:g1:plants-around-us:plant}{tumbuhan}
induknya, lewat pembelahan; muridnya bermula sebagai satu \omterm{def:g4:animal-reproduction:sexual}{sel telur}
yang terbuahi; tiap rantai yogurtnya bermula sebagai satu \omterm{def:g6:producing-our-food:fermentation}{mikroba} yang
membelah lalu membelah lagi. Segala yang digambar berpangkal pada satu
\omterm{def:g5:first-look-at-cells:cell}{sel} awal.

\textbf{12.} Kalimat 1 --- setiap \omterm{def:g6:exploring-our-environment:components}{makhluk hidup} tersusun atas \omterm{def:g5:first-look-at-cells:cell}{sel}:
paling baik diperlihatkan A dan B (dan oleh keseluruhan kelilingnya).
Kalimat 2 --- setiap \omterm{def:g5:first-look-at-cells:cell}{sel} dari sebuah \omterm{def:g5:first-look-at-cells:cell}{sel} yang membelah: paling baik
diperlihatkan C dan D, tempat pembelahan menjadi cara hidup yang
terlihat.
\end{solution}
